%==============================================================================
%  CAPÍTULO — ARQUITECTURA DE LA LEY
%==============================================================================
\chapter{Arquitectura de la Ley \Nro 20.720: mapa de procedimientos e instituciones}
\label{cap:arquitectura}

\fichaCapitulo{Visión de conjunto del sistema.}{Estructura por capítulos de la ley,
catálogo completo de procedimientos concursales, órganos e instituciones, y ubicación de
cada materia dentro de la obra.}

%==============================================================================
\section{Por qué un mapa previo}
%==============================================================================

La \LIR{} \citeyearpar{ley-20720} no se ordena por instituciones, sino por sujetos y
procedimientos. Quien la consulta buscando una figura ---la verificación de créditos, la
junta de acreedores, la continuidad de giro--- descubre que no está regulada en un solo
lugar, sino repetida, con variantes, en cada uno de los procedimientos que la emplean.
A ello se añadió la \Lreforma{} \citeyearpar{ley-21563}, que insertó procedimientos
simplificados mediante artículos con numeración literal ---273 A, 286 S--- sin alterar la
numeración preexistente.

El resultado es un cuerpo normativo cuya lectura lineal resulta desorientadora. Este
capítulo ofrece el mapa que permite recorrerlo: la estructura formal de la ley, el
catálogo íntegro de sus procedimientos e instituciones, y la remisión al lugar de esta
obra donde cada materia se desarrolla.

%==============================================================================
\section{Estructura formal de la ley}
%==============================================================================

\begin{table}[htbp]
\centering
\small
\caption{Capítulos de la Ley \Nro 20.720 y su contenido}
\label{tab:estructura-ley}
\begin{tabularx}{\textwidth}{@{}p{0.07\textwidth}>{\raggedright\arraybackslash}p{0.30\textwidth}X@{}}
\toprule
\textbf{Cap.} & \textbf{Denominación} & \textbf{Contenido} \\
\midrule
\textsc{i} & Disposiciones generales
& Ámbito de aplicación, definiciones, competencia, régimen de recursos, notificaciones y
Boletín Concursal \\
\addlinespace
\textsc{ii} & De los veedores y liquidadores
& Nóminas, nominación, deberes, inhabilidades, honorarios, cuentas y régimen de
responsabilidad de los auxiliares \\
\addlinespace
\textsc{iii} & Del procedimiento concursal de reorganización
& Reorganización judicial de la Empresa Deudora, protección financiera concursal,
acuerdos y su impugnación; acuerdos de reorganización extrajudicial o simplificados \\
\addlinespace
\textsc{iv} & Del procedimiento concursal de liquidación
& Liquidación voluntaria y forzosa de la Empresa Deudora, desasimiento, incautación,
juntas, realización del activo, pagos y sobreseimiento \\
\addlinespace
\textsc{v} & De los procedimientos concursales de la Persona Deudora
& Renegociación administrativa ante la \superir{} y liquidación de los bienes de la
Persona Deudora \\
\addlinespace
\textsc{vi} & De las acciones revocatorias concursales
& Revocación objetiva y subjetiva, período sospechoso, legitimación y plazos \\
\addlinespace
\textsc{vii} & Del arbitraje concursal
& Sometimiento a arbitraje de los procedimientos de reorganización y liquidación,
nómina de árbitros y reglas de tramitación \\
\addlinespace
\textsc{viii} & De la insolvencia transfronteriza
& Adopción de la Ley Modelo de la CNUDMI \parencite{ley-modelo-uncitral}: reconocimiento,
cooperación, COMI y medidas provisionales \\
\addlinespace
\textsc{ix} & De la Superintendencia de Insolvencia y Reemprendimiento
& Naturaleza, atribuciones normativas, fiscalizadoras y sancionatorias del órgano \\
\bottomrule
\end{tabularx}
\end{table}

\begin{alerta}[Sobre la numeración tras la reforma]
La \Lreforma{} incorporó los procedimientos simplificados sin renumerar el articulado
preexistente, mediante artículos identificados con letra (por ejemplo, 169 A, 273 A,
286 S). Las remisiones de esta obra a esos preceptos se han cotejado contra el texto refundido
vigente; con todo, al citar jurisprudencia anterior a agosto de 2023 el lector debe verificar si
el precepto invocado en el fallo conserva la numeración y el contenido actuales, pues la reforma
desplazó o redefinió varias disposiciones.
\end{alerta}

%==============================================================================
\section{Catálogo de procedimientos concursales}
%==============================================================================

El sistema chileno contempla, tras la reforma, seis vías procedimentales, más el
arbitraje como modalidad de tramitación y el reconocimiento transfronterizo como
procedimiento auxiliar.

\begin{table}[htbp]
\centering
\small
\caption{Procedimientos concursales vigentes y su tratamiento en esta obra}
\label{tab:catalogo-procedimientos}
\begin{tabularx}{\textwidth}{@{}>{\raggedright\arraybackslash}p{0.27\textwidth}
                                 >{\raggedright\arraybackslash}p{0.20\textwidth}
                                 X
                                 p{0.13\textwidth}@{}}
\toprule
\textbf{Procedimiento} & \textbf{Sujeto} & \textbf{Sede} & \textbf{Capítulos} \\
\midrule
Reorganización judicial ordinaria & Empresa Deudora mediana y grande & Juzgado civil
& \ref{cap:reorganizacion-ordinaria}, \ref{cap:tram-reorg-ordinaria} \\
\addlinespace
Reorganización simplificada & Micro y pequeña empresa & Juzgado civil
& \ref{cap:reorganizacion-simplificada}, \ref{cap:tram-reorg-simplificada} \\
\addlinespace
Acuerdo de reorganización extrajudicial & Empresa Deudora & Extrajudicial, con
aprobación judicial & \ref{cap:reorganizacion-ordinaria} \\
\addlinespace
Renegociación administrativa & Persona Deudora & \superir{}
& \ref{cap:renegociacion}, \ref{cap:tram-renegociacion} \\
\addlinespace
Liquidación ordinaria (voluntaria y forzosa) & Empresa Deudora mediana y grande
& Juzgado civil & \ref{cap:liquidacion-ordinaria}, \ref{cap:tram-liq-ordinaria} \\
\addlinespace
Liquidación simplificada & Persona Deudora y MIPE & Juzgado civil
& \ref{cap:liquidacion-simplificada}, \ref{cap:tram-liq-simplificada} \\
\addlinespace
Arbitraje concursal & Empresa Deudora & Tribunal arbitral & \ref{cap:arbitraje} \\
\addlinespace
Reconocimiento de procedimiento extranjero & Deudor con activos en Chile & Juzgado civil
& \ref{cap:transfronteriza} \\
\bottomrule
\end{tabularx}
\end{table}

%==============================================================================
\section{Catálogo de instituciones}
%==============================================================================

Las instituciones que atraviesan varios procedimientos ---y que por ello resultan
difíciles de localizar en la ley--- se sistematizan a continuación, con remisión al lugar
de la obra en que se estudian.

\subsection{Órganos y sujetos del concurso}

\begin{description}[leftmargin=0pt,style=nextline,font=\sffamily\bfseries]
  \item[Superintendencia de Insolvencia y Reemprendimiento] Órgano fiscalizador con
  potestad normativa, de registro y sancionatoria. Capítulo \ref{cap:institucionalidad};
  su producción normativa, capítulo \ref{cap:normativa-superir}.

  \item[Veedor] Auxiliar de la reorganización. Informa a los acreedores sobre la
  propuesta y fiscaliza al deudor durante la protección financiera. Capítulo
  \ref{cap:institucionalidad}.

  \item[Liquidador] Auxiliar de la liquidación. Asume la administración del patrimonio
  desasido, incauta, realiza y reparte. Capítulo \ref{cap:institucionalidad}.

  \item[Martillero concursal] Interviene en la realización del activo, con excepciones en
  la liquidación simplificada. Capítulos \ref{cap:institucionalidad} y
  \ref{cap:liquidacion-simplificada}.

  \item[Árbitro concursal] Capítulo \ref{cap:arbitraje}.

  \item[Junta de acreedores] Órgano colectivo de deliberación y decisión, en sus
  modalidades constitutiva, ordinaria y extraordinaria. Capítulos
  \ref{cap:reorganizacion-ordinaria} y \ref{cap:liquidacion-ordinaria}.

  \item[Interventor] Supervisa el cumplimiento del acuerdo de reorganización. Capítulo
  \ref{cap:reorganizacion-simplificada}.
\end{description}

\subsection{Instituciones transversales del procedimiento}

\begin{description}[leftmargin=0pt,style=nextline,font=\sffamily\bfseries]
  \item[Boletín Concursal] Medio de publicidad y notificación de todas las resoluciones.
  Capítulo \ref{cap:arquitectura}, sección \ref{sec:boletin}.

  \item[Verificación de créditos] Acto por el cual el acreedor comparece y acredita su
  crédito y preferencia. Capítulos \ref{cap:reorganizacion-ordinaria} y
  \ref{cap:liquidacion-ordinaria}.

  \item[Protección Financiera Concursal] Capítulos \ref{cap:reorganizacion-ordinaria} y
  \ref{cap:reorganizacion-simplificada}.

  \item[Desasimiento] Capítulo \ref{cap:liquidacion-ordinaria}.

  \item[Continuidad de las actividades económicas] Capítulo
  \ref{cap:liquidacion-ordinaria}.

  \item[Prelación de créditos] Aplicación concursal del \artcc{2472}
  \citeyearpar{codigo-civil}, con las preferencias laborales y previsionales. Capítulos
  \ref{cap:laboral-penal} y \ref{cap:creditos-laborales}.

  \item[Acciones revocatorias concursales] Capítulo \ref{cap:revocatorias}.

  \item[Extinción de los saldos insolutos] Capítulo \ref{cap:discharge-cae}.

  \item[Incidente de mala fe] Capítulo \ref{cap:mala-fe}.

  \item[Sobreseimiento definitivo y temporal] Capítulos \ref{cap:liquidacion-ordinaria} y
  \ref{cap:liquidacion-simplificada}.

  \item[Cuenta final de administración] Capítulos \ref{cap:institucionalidad} y
  \ref{cap:liquidacion-ordinaria}.
\end{description}

%==============================================================================
\section{El Boletín Concursal como eje del sistema de notificaciones}
\label{sec:boletin}
%==============================================================================

Ninguna institución de la ley condiciona tanto la práctica cotidiana como el
\boletin{}. La regla general es que las resoluciones se notifican mediante su
publicación, sin necesidad de actuación del receptor ni de comparecencia de las partes.

Las consecuencias operativas son tres, y las tres explican una porción significativa de
los perjuicios profesionales en materia concursal:

\begin{enumerate}
  \item \textbf{Los plazos corren sin aviso.} No hay notificación personal que active la
  atención del abogado. La revisión diaria del \boletin{} es la única forma de no perder
  un plazo fatal.

  \item \textbf{La publicidad es total.} Cualquiera puede consultar la situación
  concursal del deudor. El cliente debe saberlo antes de presentar.

  \item \textbf{La carga de la publicación recae en el auxiliar.} Los deberes de
  publicación del veedor y del liquidador, y el régimen de notificaciones por correo
  electrónico, están regulados por la \ncg{14} \citeyearpar{ncg-14}.
\end{enumerate}

\begin{foco}[Rutina mínima de control]
Revisión diaria del \boletin{} por RUT del deudor, con responsable identificado dentro del
estudio y registro de la revisión. Es el control más barato y el que evita el mayor
número de contingencias.
\end{foco}

El fundamento normativo está en el \art{6}: cuando el tribunal ordena notificar por avisos, la
notificación se practica mediante \destacar{publicación en el \boletin{}}, y se entiende hecha
\destacar{desde la fecha de su inserción}. La publican el veedor, el liquidador o la \superir{}
dentro de los dos días siguientes a la resolución.

%==============================================================================
\section{Disposiciones generales: competencia, plazos, incidentes y quórums}
%==============================================================================

Cuatro reglas del Capítulo I de la ley gobiernan la mecánica de todos los procedimientos y
conviene fijarlas antes de entrar a cada uno.

\subsection{Competencia}

El \art{3} radica los procedimientos concursales en el \destacar{juzgado de letras del domicilio
del deudor}. El acreedor puede promover el incidente de incompetencia conforme a las reglas
generales. En las ciudades asiento de Corte, la distribución de causas se rige por un auto
acordado de la Corte de Apelaciones respectiva, con radicación preferente de las causas
concursales.

\subsection{Cómputo de plazos}

El \art{7} establece la regla que ordena todo el calendario procesal: los plazos de días son de
\destacar{días hábiles} ---inhábiles los domingos y festivos--- salvo que la ley diga
expresamente que son de días corridos. Se cuentan desde el día siguiente a la notificación. Y
los plazos para actuar \emph{antes} de una fecha se cuentan \destacar{hacia atrás} desde el día
anterior a la actuación.

\begin{alerta}[Hábiles, corridos y hábiles administrativos: tres cómputos distintos]
La ley usa tres relojes. Los plazos judiciales, por regla general, son de días hábiles
(art. 7). Algunos se declaran de días corridos ---por ejemplo, los noventa días de la
renegociación (\art{260})---. Y los plazos administrativos ante la \superir{} se cuentan en días
hábiles administrativos (\art{262}). Confundirlos es la fuente de error de plazo más común del
sistema.
\end{alerta}

\subsection{Incidentes}

El \art{5} restringe la incidencia: \destacar{sólo pueden promoverse incidentes en las materias
que la ley permite expresamente}. Se tramitan conforme al \cpc{} y \destacar{no suspenden el
procedimiento}, salvo que la ley disponga lo contrario. Es una manifestación más del principio
de celeridad: la ley no tolera que la incidencia paralice el concurso.

\subsection{Los quórums del artículo 2}

Las decisiones de la junta se adoptan por tres quórums que la ley define y que se invocan a lo
largo de toda la obra:

\begin{description}[leftmargin=0pt,style=nextline,font=\sffamily\bfseries]
  \item[Quórum Simple] La mayoría del pasivo verificado con derecho a voto presente en la junta.
  \item[Quórum Calificado] La mayoría absoluta del pasivo con derecho a voto.
  \item[Quórum Especial] Los \destacar{dos tercios} del pasivo total con derecho a voto.
\end{description}

\subsection{Exigibilidad y el principio de especialidad}

El \art{8} zanja la relación de la ley concursal con el resto del ordenamiento: \destacar{las
normas de leyes especiales prevalecen} sobre las de esta ley, y lo no regulado por aquéllas se
rige supletoriamente por ésta.

\begin{foco}[El artículo 8 es el eje del conflicto del CAE]
Esta regla de prevalencia de la ley especial es, precisamente, el fundamento que la
jurisprudencia de exclusión invoca para sustraer el Crédito con Aval del Estado del
\discharge{}: si la legislación de educación superior es «especial», prevalecería sobre la
extinción concursal. La discusión ---y su reverso, la \lat{lex posterior}--- se desarrolla en el
capítulo \ref{cap:discharge-cae}.
\end{foco}
